\chapter{Detalles de Implementaci\'on}
Para este momento, ya se tienen definidos los requerimientos funcionales y la
arquitectura de software que mejor se adec\'ua a la aplicaci\'on DermaUH 3.0. 
En este cap\'itulo se abordar\'an cuestiones relacionadas concretamente con la 
implementaci\'on de la aplicaci\'on.

\section{Stack Tecnol\'ogico}
Como se ha mencionado a lo largo de este documento, DermaUH 3.0 es una aplicaci\'on
web, la cual tendr\'a integrado un modelo de Aprendizaje de M\'aquinas (ML por sus
siglas en ingl\'es). Por tal motivo se decidi\'o optar por un conjunto de tecnolog\'ias
ideales para el desarrollo web y que se puedan combinar sin problemas con t\'ecnicas de ML.
En el cap\'itulo anterior se seleccion\'o la arquitectura por capas como patr\'on principal
dise\~no, dividiendo la aplicaci\'on en tres niveles: Presentaci\'on (UI, frontend), L\'ogica
y Reglas de Negocio (backend) y Base de Datos (BD). Cada nivel tiene implementado con 
su propia tecnolog\'ia, seleccionada luego de analizar los requerimientos de la aplicaci\'on,
la experiencia del desarrollador y el tiempo disponible para la implementaci\'on. 

\subsection*{React}
\begin{figure}[ht] 
    \centering
    \includegraphics[width=0.2\textwidth]{Graphics/react.png}
    \caption{Logo de React} 
    \label{React}
\end{figure}
Para la capa de presentación se ha seleccionado React, una biblioteca de JavaScript 
desarrollada por Meta para la construcción de interfaces de usuario basadas en componentes 
reactivos \cite{react2021}. React ofrece ventajas determinantes para DermaUH 3.0, como el 
renderizado eficiente mediante el DOM virtual, que garantiza tiempos de respuesta rápidos 
incluso en dispositivos con recursos limitados, y la capacidad de crear una interfaz 
responsive que se adapta a resoluciones de 320px a 1920px, tal como exige el requisito no 
funcional RNF4. Su ecosistema permite integrar fácilmente gestores de estado como Redux o 
contextos nativos, lo que facilita la gestión de las sesiones de usuario autenticado con 
JWT y la visualización de historiales clínicos, lesiones y diagnósticos, adaptándose 
perfectamente a los casos de uso identificados para los actores Doctor y Administrador.

\subsection*{Django}
\begin{figure}[ht] 
    \centering
    \includegraphics[width=0.2\textwidth]{Graphics/django.png}
    \caption{Logo de Django} 
    \label{Django}
\end{figure}
Para la capa de lógica de negocio se ha elegido Django, un framework de desarrollo web en 
Python que sigue el patrón MTV (Model-Template-View) y que incluye un potente sistema de 
rutas, autenticación y un ORM integrado \cite{django2022}. Django proporciona ventajas 
críticas para DermaUH 3.0, como su módulo de autenticación que permite implementar de 
forma nativa el control de acceso basado en roles (administrador, doctor, residente) y la 
generación de JSON Web Tokens para cumplir con RNF1. Además, su arquitectura basada en 
aplicaciones reutilizables facilita la organización de los servicios de gestión de pacientes, 
lesiones, diagnósticos y la comunicación con el modelo de inteligencia artificial, permitiendo 
cumplir con los requisitos funcionales RF1 a RF5 y mantener un tiempo de respuesta 
inferior a dos segundos (RNF3) mediante el uso de vistas asíncronas y caché.

\subsection*{PostgreSQL}
\begin{figure}[ht] 
    \centering
    \includegraphics[width=0.2\textwidth]{Graphics/postgres.png}
    \caption{Logo de PostgreSQL} 
    \label{PostgreSQL}
\end{figure}
Para la capa de persistencia se ha optado por PostgreSQL, un sistema gestor de bases de 
datos relacional de código abierto reconocido por su solidez, cumplimiento de estándares 
SQL y soporte de tipado avanzado \cite{postgresql2023}. PostgreSQL se adapta a DermaUH 3.0 
porque permite implementar exactamente el modelo entidad-relación definido en la figura 
\ref{MER}, incluyendo claves foráneas con reglas ON DELETE CASCADE para mantener la 
integridad referencial entre pacientes, lesiones e imágenes, así como campos específicos 
como fototipos (enteros acotados) y rutas de archivos. Su soporte para conexiones cifradas 
mediante SSL/TLS garantiza el cumplimiento de RNF2, y su motor de consultas optimizado 
asegura que las operaciones de listado, búsqueda y filtrado de pacientes 
(con índices adecuados) respondan en los márgenes exigidos por RNF3 incluso con 50 
usuarios concurrentes.

\section{Implementaciones}
En el apartado anterior, se definieron las tecnolog\'ias utilizadas para el desarrollo de la 
aplicaci\'on web DermaUH 3.0. En este apartado se explicar\'a de cara de al desarrollador, la
implementaci\'on concreta utilizada en cada capa, atendiendo al stack seleccionado en el apartado
anterior. Se mostrar\'an adem\'as ejemplos de c\'odigo y como se resolvieron ``problemas''
que fueron surgiendo durante el proceso de desarrollo.

\subsection{Backend}
Como se mencion\'o en el cap\'itulo anterior, la arquitectura por capas posee un alto dinamismo,
por tanto se pueden adoptar enfoques h\'ibridos y cada nivel puede tener su propia implementaci\'on.
Para organizar el backend, se optó por dividir la lógica en varias aplicaciones Django 
independientes (auth\_service, email\_service, medical\_service, image\_service). 
Cada una se encarga de un dominio específico, lo que facilita el mantenimiento y la 
evolución por separado. Aunque todas comparten la misma base de datos y se ejecutan en 
el mismo servidor (por lo que no se trata de una arquitectura de microservicios completa),
esta modularidad permite, por ejemplo, reemplazar o actualizar el servicio de imágenes 
sin afectar al resto, siempre que se respeten las interfaces de comunicación. La Figura~\ref{microserviciosDjango} 
muestra como fue dividido el backend en aplicaciones m\'as peque\~nas y de responsabilidad \'unica.
Se opt\'o por este enfoce modular por cuestiones de optimizaciones en las peticiones,
ahora en vez de existir una sola aplicaci\'on que consuma todas las solicitudes, existen varias 
que leen y responden a solicitudes m\'as concretas.

\begin{figure}[ht] 
    \centering
    \includegraphics[width=0.4\textwidth]{Graphics/microservicios.png}
    \caption{Servicios Django, la carpeta logs es para almacenar informaci\'on de cada servicio} 
    \label{microserviciosDjango}
\end{figure}

A continuaci\'on una breve explicaci\'on de cada servicio y que tareas cumple:
\begin{itemize}
    \item \textbf{auth\_service}: Las tareas de este servicio son las relacionadas con registro e inicio de sesi\'on. Adem\'as maneja todas las operaciones CRUD sobre los usarios y es donde se aceptan o rechazan las solicitudes. Tambi\'en genera el JWT correspondiente una vez que un usuario incie sesi\'on en el sistema.
    \item \textbf{email\_service}: Es donde se almacena el cliente de correo y donde se generan los gmails que reciben los usuarios en relaci\'on a su solicitud.
    \item \textbf{medical\_service}: Controla toda la gesti\'on de pacientes, lesiones y diagn\'osticos. Es el servicio que procesa las solicitudes provenientes de usuarios con el rol de Doctor.
    \item \textbf{image\_service}: Controla la parte de la aplicaci\'on dedicada a las imagenes: vincularlas a las lesiones y subirlas al sistema. En este servicio tambi\'en se manejan las peticiones al modelo de ML vinculado a la aplicacion.
\end{itemize}

Antes de comenzar el an\'alisis detallado de cada servicio, es v\'alido esclarecer
conceptos que son muy habituales en el desarrollo web con Django como framework. Se 
incluir\'an tambi\'en im\'agenes de c\'odigo para un mejor entendimiento al lector.
Toda la informaci\'on se extrajo de la \href{https://docs.djangoproject.com/en/6.0/}
{Documentaci\'on oficial de Django} y del foro \href{https://www.django-rest-framework.org/}{Django Rest-Framework}\footnote{Django RestFramework es una extensi\'on de Django muy utilizada para la construcci\'on de Web APIs a base de Django}.

\subsubsection*{Model}
Seg\'un la documentaci\'on oficial de Django\cite{django2022}, los Models son clases 
que se utilizan para mapear tablas de la base de datos. En s\'intesis, son la representaci\'on
en c\'odigo de las entidades. Si una entidad posee una propiedad de tipo entero, entonces
el modelo en Django, tendr\'a un campo de tipo IntegerField\footnote{Campo utilizado para representar enteros}.

La Figura~\ref{modelDermatologist} muestra un ejemplo de como definir una entidad en Django,
concretamente la entidad derm\'atologo. En la figura se muestra que la clase \texttt{Dermatologist}
hereda de AbstractUser y PermissionMixin\footnote{Estas clases definen que el modelo en cuesti\'on la aplicaci\'on debe tratarlo como usuarios y que puede tener uno o varios permisos respectivamente}
ambas heredan de la clase Model de Django, por lo tanto Dermatologist hereda de Model.
\begin{figure}[ht] 
    \centering
    \includegraphics[width=0.6\textwidth]{Graphics/modelDermatologist.png}
    \caption{Modelo Dermatolgist, implementado en Django} 
    \label{modelDermatologist}
\end{figure}

\subsubsection*{Serializer}
Los serializers (serializadores) en Django, funcionan como un traductor de informaci\'on. Permiten
que datos complejos como consultas e instancias de modelos, se conviertan en datos nativos de Python.
La forma en que se utilizaron en el desarrollo de DermaUH 3.0 fue para manejar las peticiones
de tipo POST\footnote{Peticiones de crear nuevas instancias de modelos para la BD}, PUT\footnote{Peticiones para actualizar valores de instancias guardadas en BD}
y como DTO\footnote{Data Transfer Object, como su nombre indican, se utilizan para extraer informaci\'on de la BD sin necesidad de devolver toda la instancia\cite{DTO}}.
La Figura~\ref{serializerDTO} evidencia el uso como DTO de un serializer, mientras que la
Figura~\ref{serializerPost} muestra un ejemplo de un serializer usado para ``atender'' una solicitud
de tipo POST.

\begin{figure}[ht] 
    \centering
    \includegraphics[width=0.6\textwidth]{Graphics/dtoSerializer.png}
    \caption{Serializer utilizado como DTO} 
    \label{serializerDTO}
\end{figure}

\begin{figure}[ht] 
    \centering
    \includegraphics[width=0.6\textwidth]{Graphics/postSerializer.png}
    \caption{Serializer para crear una solicitud de dermat\'ologo} 
    \label{serializerPost}
\end{figure}

Como se muestra adem\'as en la figura anterior, con el serializer se puede hacer validaci\'on
de datos (funci\'on validate) y adem\'as se pueden crear y guardar instancias. Esto es muy ventajoso
en el desarrollo de aplicaciones, los serializers tiene la capacidad de traducir datos, validarlos y crear o actualizar instancias.

\subsection*{View}
Las Views (vistas) en Django son el portal de la comunicaci\'on con aplicaciones externas, como el frontend.
Su principal funci\'on es recibir las solicitudes y luego tomar la decisi\'on de consular un serializer o no. 
Generalmente las views que reciben peticiones de tipo GET o DELETE, no requieren de un serializer,
tiene sentido ya no tienen que actualizar ni escribir en BD. Django aporta la clase \texttt{APIView}
que es la m\'as utilizada en el desarrollo, sin embargo existen otras como ListAPIView o DestroyAPIView
que tambi\'en pueden usarse para listar o eliminar elementos respectivamente. La Figura~\ref{postView}
muestra una vista que maneja una petici\'on POST, n\'otese como en el apartado ``request'' (cuerpo de la petici\'on)
se hace referencia al serializer mostrado en la Figura~\ref{serializerPost}, evidenciando
que el cuerpo de la solicitud en cuesti\'on debe ser igual al de la clase \texttt{Meta} en el serializer.

\begin{figure}[ht]
    \centering
    \includegraphics[width=0.6\textwidth]{Graphics/postView.png}
    \caption{View para peticiones Post} 
    \label{postView}
\end{figure}

La Figura~\ref{listView} muestra una view para peticiones GET, concretamente, el listado con
filtro de los dermat\'ologos (usuarios) del sistema. En esta view se definen los diferentes tipos
de filtros, se prepara el conjunto de datos a los cuales se les har\'a la consulta y finalmente s
se devuelven los usuarios correspondientes. Todo ese proceso lo realiza la clase \texttt{ListAPIView},
su funcionamiento puede ser consultado en la documentaci\'on oficial de RestFramework.\newpage

\begin{figure}[ht] 
    \centering
    \includegraphics[width=0.6\textwidth]{Graphics/getView.png}
    \caption{View para listado de usuarios} 
    \label{listView}
\end{figure}

Existe adem\'as una alternativa de Views, que es la ModelViewSet. Este objeto proporciona
un conjunto de vistas predefinidas, que cumplen con las operaciones CRUD. Esta opci\'on
es totalmente v\'alida y utilizable en proyectos de programaci\'on. Sin embargo para el 
desarrollo de DermaUH 3.0 siempre se opt\'o por un enfoque modular, en todos los sentidos,
el ModelViewSet agrupa todas las View en un solo objeto, si por alg\'un motivo desconocido
una de las vistas deja de responder, esto puede provocar que el resto de las operaciones que
engloba el conjunto de Views, sean inaccesibles hasta que se solucione el error. Al separarlos
en varias vistas, permite que si alguna est\'a desconectada, o en mantenimiento, los usuarios
puedan seguir utilizando otras funcionalidades sin problema.

Una vez entendido los conceptos claves de la programaci\'on con Django, a continuaci\'on
se har\'a una explicaci\'on m\'as detallada de cada uno de los servicios que componen el
backend de la aplicaci\'on. 

\subsubsection{Servicio de Autenticaci\'on (auth\_service)}
Como se mencionó al inicio de este apartado, las tareas del auth\_service son aquellas 
relacionadas con la autenticación, registro y solicitudes de los usuarios del sistema, 
incluyendo la generación del JWT tras un inicio de sesión exitoso. Para garantizar la 
seguridad de las credenciales, las contraseñas se almacenan aplicando el algoritmo bcrypt 
con un factor de costo 12, nunca en texto plano. El JWT generado incluye los claims 
user\_id, role (admin, doctor o residente) y un tiempo de expiración de 8 horas; las 
sesiones prolongadas se soportan mediante un token de refresco almacenado en una cookie 
HttpOnly, con renovación automática. El servicio también valida las solicitudes de 
registro comprobando la unicidad del email y del carnet de identidad, persistiendo la 
solicitud en estado pendiente y notificando al administrador por correo. La Figura~\ref{tokenGen}
muestra como se genera un token personalizado para poder establecer comunicaci\'on
entre los distintos servicios.

\begin{figure}[ht] 
    \centering
    \includegraphics[width=0.6\textwidth]{Graphics/token.png}
    \caption{Generaci\'on del token} 
    \label{tokenGen}
\end{figure}

\subsubsection{Servicio de Correo Electrónico (email\_service)}
El servicio de correo electrónico se ocupa de la notificación asíncrona a los actores del 
sistema, sin requerir autenticación previa para recibir las peticiones. Este servicio 
gestiona tres tipos de envíos. En primer lugar, cuando un nuevo dermatólogo o residente 
completa el formulario de registro, el servicio notifica al administrador con los datos 
básicos del solicitante; si se trata de un residente, además se envía una notificación al 
tutor indicado para que conozca la petición. En segundo lugar, tras la revisión del 
administrador, el servicio informa al usuario si su solicitud ha sido aceptada o rechazada,
adaptando el mensaje al resultado. Por último, cuando un usuario olvida su contraseña, 
el servicio genera y remite un enlace seguro para restablecerla. Todos los 
correos se entregan mediante la función \texttt{send\_mail} de Django, y la configuración 
del remitente se obtiene de las variables de entorno, garantizando así la separación 
entre configuración y código.

\subsubsection{Servicio Clínico (medical\_service)}
El servicio clínico constituye el núcleo funcional de la aplicación para los dermatólogos, 
ya que encapsula todas las operaciones relacionadas con la gestión de pacientes, 
lesiones y diagnósticos. A través de este servicio, el doctor autenticado puede 
registrar nuevos pacientes, consultar sus datos personales y de contacto, filtrar la 
lista de pacientes mediante búsquedas y eliminar aquellos que ya no estén bajo su 
seguimiento, dando cumplimiento al requisito funcional RF1. Sobre las lesiones, el 
servicio permite registrar su localización anatómica y fototipo cutáneo, 
y almacenar los diagnósticos definitivos emitidos por el médico, manteniendo siempre la 
integridad referencial mediante las reglas de eliminación en cascada definidas en el 
modelo de datos. Además, implementa la funcionalidad de compartir lesiones con otros 
doctores, lo que permite que un especialista autorizado acceda a la información 
clínica de un paciente que no está directamente bajo su cargo, cumpliendo así el 
requisito RF2.

\subsubsection{Servicio de Imágenes e Inteligencia Artificial (image\_service)}
El servicio de imágenes se especializa en la recepción, validación y procesamiento de 
las fotografías clínicas y dermatoscópicas enviadas por los dermatólogos. Una vez que el 
usuario autenticado sube una imagen asociada a una lesión, este servicio se encarga de 
almacenar los metadatos correspondientes (ruta del archivo, tipo de imagen, comentarios) 
en la base de datos y de gestionar el archivo físico en el servidor. Su funcionalidad más 
relevante consiste en actuar como puente hacia el modelo de inteligencia artificial: 
cuando el médico solicita una sugerencia de clasificación, el image\_service envía la 
imagen al modelo previamente entrenado, recibe el resultado (tipo de lesión y nivel de 
confianza) y lo persiste como un diagnóstico automático pendiente de validación. 
De esta forma se satisface el requisito funcional RF3, garantizando que el especialista 
reciba una ayuda inmediata sin que el sistema reemplace su criterio clínico, 
y todo ello respetando el tiempo máximo de respuesta exigido en RNF3.

\subsubsection{Comunicaci\'on entre Servicios}
Cada servicio se comporta como una WebApi implementada con Django. Aunque no exista
una dependencia directa entre estos, durante la etapa de desarrollo, surgieron inconvenientes
que obligaron establecer de alguna forma comunicaci\'on entre los servicios. Un ejemplo 
de ello es la creaci\'on de un paciente, si se observa la descripci\'on de esta entidad
(P\'ag 37-38), se observa que tiene una llave for\'anea referenciando a la entidad Dermat\'ologo.
Sin embargo, el servicio medical\_service (donde se registran pacientes) no tiene conocimiento
de la existencia de esta entidad. La soluci\'on est\'a dada por el JWT generado en el auth\_service,
como se explic\'o anteriormente: el token generado guarda el claim del user\_id, que es el
identificador en base de datos del usuario actual. Por tanto cuando el usuario actual registre
un paciente, el valor de la llave for\'anea a la entidad dermat\'ologo ser\'a ese user\_id
proveniente del token generado. La Figura~\ref{postPatient} muestra como el user\_id proveniente
de la solicitud es enviado al serializer.

\begin{figure}[ht] 
    \centering
    \includegraphics[width=0.6\textwidth]{Graphics/postPatient.png}
    \caption{Uso de user\_id proveniente del JWT} 
    \label{postPatient}
\end{figure}

\subsection{Frontend}
La capa de presentaci\'on, fue desarrollada con el framework React. Para ello se opt\'o
por un patr\'on propio del framework: La Arquitectura basada en Caracter\'isticas (Feature Architecture)\cite{reactArchitecture}
que consiste en dividir la aplicaci\'on frontend, en subaplicaciones m\'as peque\~nas, permitiendo
una mayor capacidad de escalabilidad. Para el caso de DermaUH cada feature
(caracter\'stica) es una de las entidades principales del proyecto (Paciente, Dermat\'ologos, Lesiones).
Este enfoque permite, al igual que en el backend, una separaci\'on de 
responsabilidades, donde cada feature se encarga de tareas \'unicas y no dependen una
de la otra para poder ejecutarse. Con esta idea se logra que si por alg\'un motivo, por ejemplo,
el registro de pacientes deja de responder desde el backend, la aplicaci\'on contin\'ue
funcionando en el resto de los aspectos. Al observar la Figura~\ref{arquitecturaReact}
se puede apreciar la distribuci\'on de archivos de la arquitectura basada en caracter\'isticas.\newpage

\begin{figure}[ht] 
    \centering
    \includegraphics[width=0.4\textwidth]{Graphics/features.png}
    \caption{Arquitectura basada en caracter\'isticas en el frontend DermaUH} 
    \label{arquitecturaReact}
\end{figure}

Al igual que en el apartado Backend, se realizar\'a un recorrido por diferentes conceptos
que son importantes para entender las ideas detr\'as de la implementaci\'on del frontend
con React.

\subsubsection*{Axios}
Axios es una biblioteca de JavaScript que ayuda a la realizaci\'on de peticiones HTTP desde
la aplicaci\'on actual, hacia otros servidores o aplicaciones\cite{axios}. Se utiliza definiendo la 
direcci\'on (URL) del sitio al cual se realizar\'an las peticiones. En caso de DermaUH, con backend
en Django, este genera una URL para cada uno de los servicios\footnote{Recordar que se comportan como Web APIs independientes},
entonces desde la capa de presentaci\'on, se define un axios para cada servicio del backend.
La Figura~\ref{axiosFigure} muestra como se usa axios para iniciar la comunicaci\'on con el medical\_service.
Haciendo uso del token guardado en el localStorage\footnote{Es un diccionario global en el cual se pueden almacenar todo tipo de valores, en este caso se almacena el access token devuelto por el backend y se consume en el axios para enviar autenticaci\'on al resto de servicios}
permite autenticarse en dicho servicio y utilizarlo si el usuario cuenta con los permisos correspondientes.
Destacar tambi\'en que axios establece comunicaci\'on de forma as\'incrona mediante Promesas (Promise)\cite{promise}.

\begin{figure}[ht] 
    \centering
    \includegraphics[width=0.5\textwidth]{Graphics/axios.png}
    \caption{Configuraci\'on Axios para establecer comunicaci\'on} 
    \label{axiosFigure}
\end{figure}

\subsubsection*{Services}
Los Services (servicios) en React son funciones de corta extensión, normalmente no superior
a las diez líneas de código, cuya responsabilidad exclusiva es realizar peticiones al 
backend utilizando el mensajero Axios previamente configurado. Actúan como una capa de 
abstracción que aísla la lógica de comunicación HTTP del resto de la aplicación, de modo 
que los componentes y hooks no necesitan conocer los detalles de las URLs, los métodos 
HTTP ni los encabezados de autenticación. En DermaUH 3.0 se ha organizado un servicio por 
cada entidad principal (Paciente, Lesión, Diagnóstico, Autenticación) y por cada servicio 
del backend (auth\_service, medical\_service, image\_service, email\_service). 
Cada servicio exporta funciones asíncronas que, utilizando la palabra clave \texttt{await},
invocan a Axios y retornan directamente la respuesta (generalmente en formato JSON) 
sin aplicar transformaciones adicionales, manteniendo así el principio de responsabilidad 
única. 

La Figura~\ref{services} muestra como se configuran las funciones correspondientes a los servicios de la entidad Paciente,
todas ellas añaden automáticamente el token JWT almacenado en \texttt{localStorage} a la cabecera 
\texttt{Authorization}, gracias a un interceptor configurado globalmente en Axios. 
Además, estos servicios incorporan un manejo básico de errores: si la petición falla 
(por problemas de red, autenticación o validación), lanzan una excepción que será 
capturada por el hook o componente que invocó al servicio, permitiendo mostrar 
mensajes claros al usuario como exige el requisito no funcional RNF5. 
De esta forma, los servicios actúan como el puente perfecto entre la interfaz de usuario 
y la lógica de negocio remota, facilitando las pruebas unitarias (pueden ser mockeados 
fácilmente) y la evolución independiente de cada recurso.

\begin{figure}[ht] 
    \centering
    \includegraphics[width=0.6\textwidth]{Graphics/patientService.png}
    \caption{Servicios de la entidad Paciente} 
    \label{services}
\end{figure}

\subsubsection*{Hooks}
Los Hooks son funciones introducidas en React 16.8 que permiten a los componentes 
funcionales gestionar estado interno, efectos secundarios y contexto. Es una buena forma
de encapsular los servicios definidos. La Figure~\ref{patientHook} muestra como 
en un solo Hook se pueden configurar y guardar todos los servicios de la entidad 
paciente, como un diccionario en cualquier otro lenguaje de programaci\'on. Debido a su extensi\'on
solo muestran las funciones que invocan a los servicios de crear y listar pacientes.

\begin{figure}[ht] 
    \centering
    \includegraphics[width=0.6\textwidth]{Graphics/hooks.png}
    \caption{Hook para los servicios de pacientes} 
    \label{patientHook}
\end{figure}

Este hook fue implementado para encapsular toda la l\'ogica relacionada con pacientes
en un solo lugar, sin dispersarla entre m\'ultiples componentes. Esto permite reutilizar
esa misma l\'ogca en diferentes partes de la aplicaci\'on (ejemplo PatientList, PatientForm y
PatientDetails) sin duplicar c\'odigo. Además, al centralizar el manejo de estados como 
loading, error y los datos de pacientes, se facilita el mantenimiento y las pruebas 
unitarias, ya que los componentes solo se comunican con el hook sin conocer los detalles 
de implementación de las peticiones HTTP.

\subsubsection*{Components}
Los componentes son los bloques de construcción fundamentales de cualquier interfaz en 
React; cada componente es una pieza de UI independiente y reutilizable que puede recibir 
datos de entrada (props) y mantener su propio estado interno\cite{reactComponents}. 
En DermaUH 3.0 se ha adoptado una arquitectura basada en características (Feature 
Architecture), lo que significa que los componentes se agrupan por dominio funcional: 
Paciente, Dermatólogo, Lesión y Diagnóstico. Por ejemplo, dentro de la característica 
Paciente existen componentes como \texttt{PatientList} (tabla o tarjeta de pacientes con 
botones de acción), \texttt{PatientForm} (formulario de registro y edición) y 
\texttt{PatientDetails} (vista detallada con pestañas para lesiones e historial). 
Cada uno de estos componentes se comunica con los servicios definidos en la capa inferior 
(los que utilizan axios) para solicitar o enviar datos al backend, y utilizan hooks para 
manejar el estado de carga, los errores y la respuesta. La separación en componentes 
pequeños y con una única responsabilidad facilita las pruebas unitarias, el mantenimiento 
y la evolución independiente de cada funcionalidad, de modo que si en el futuro se 
desea modificar el aspecto visual de la lista de pacientes, no es necesario alterar el 
formulario de registro ni los componentes de lesiones.

\section{Modelos Machine Learning}
Como se ha hablado a lo largo de este documento, uno de los objetivos fundamentales de 
la versi\'on 3.0 de DermaUH es la incorporaci\'on de nuevos modelos de procesamiento de 
im\'agenes entrenados con un conjunto de im\'agenes cubanas. Los modelos que se utilizaron
fueron los desarrollados por el Licenciado en Ciencias de la Computaci\'on Daniel Abad, 
para m\'as referencia puede consultar su tesis de diploma\cite{tesisDaniel}. Se siguieron
rigurosamente todas las instrucciones expuestas en el documento relacionada con las transformaciones
del dataset en caso de ser necesarias.

A pesar de que utilizar pieles cubanas era un objetivo primordiar de esta nueva versi\'on,
las entidades del MINSAP no mostraron una respuesta positiva durante el per\'iodo
de recolecci\'on de im\'agenes, obteniendo solamente 50 en total. Por tal motivo
el dataset utilizado tiene un total de 4576 im\'agenes. Se extrajo de la plataforma
Kaggle\footnote{Plataforma especializada en estudios estad\'isticos y trabajos de aprendizaje de m\'aquinas},
y es una mezcla entre varios archivos dedicados a la dermatoscop\'ia. 
el dataset Fue dividido en cuatro clases (Ya tra\'ia una divisi\'on de 8 clases, se tomaron)
Melanoma, Carcinoma Basal, Carcinoma Escamoso y Otras. La distribuci\'on 
de im\'agenes por categor\'ias qued\'o de la siguiente manera:

\begin{itemize}
    \item Melanoma: 1000 para entrenamiento, 100 para pruebas y 100 para validaci\'on.
    \item Carcinoma Basal: 1000 para entrenamiento, 100 para pruebas y 100 para validaci\'on.
    \item Carcinoma Escamoso: 500 para entrenamiento, 64 para pruebas y 62 para validaci\'on.
    \item Otras lesiones: 1200 para entrenamiento, 200 para pruebas y 150 para validaci\'on.
\end{itemize}

Se expandi\'o el conjunto de Carcinoma Escamoso para hacer que la cantidad de im\'agenes por 
categor\'ias sea la misma y evitar sesgos. La t\'ecnica utilizada fue tomar la imagen original
y su versi\'on rotada 90 grados. Por tanto la cantidad de im\'agenes de Carcinoma Escamoso se 
duplic\'o, teniendo as\'i 1000 im\'agenes de esta categor\'ia. Si se analiza la tesis de Daniel\cite{tesisDaniel}
se explica que tener un dataset desequilibrado puede provocar sobreajuste y malos resultados, 
por ese motivo en esta tesis, no se hablar\'a de experimentos con el dataset sin configuraci\'on.

Al no utilizar un dataset de pieles cubanas, las m\'etricas obtenidas no son las esperadas y 
el comportamiento del modelo tampoco ser\'a el deseado. Las fuentes de las que Kaggle obtiene
conjuntos de im\'agenes, son las analizadas en el estado del arte de esta tesis (ver cap\'itulo 2).
Por tanto presentan todas las limitaciones mencionadas en su an\'alisis, incumpliendo con 
un requerimiento de la nueva versi\'on.

El modelo utilizado es un h\'ibrido dividido en dos fases de clasificaci\'on.
Una primera fase binaria (melanoma o no melanoma),
realizada con una SVM (M\'aquina de Soporte Vectorial)\cite{svm} entrenada sobre vectores
de caracter\'isticas extra\'idos con DINOv2\cite{dinov2}. La segunda fase es una 
clasificaci\'on multiclase (Carcinoma Basal, Carcinoma Escamoso, Otras lesiones), realizada 
con mediante un Vision Transformer\cite{visionTransformer} preentrenado en 
ImagenNet-21k al que se aplic\'o Fine Tunin\cite{tesisDaniel}.
El an\'alisis realizado en la tesis de Daniel\cite{tesisDaniel} determina que esta combinaci\'on
tiene una alta sensibilidad a la detecci\'on del melanoma, con un $99.7\%$, lo cual en su momento
(versiones anteriores de DermaUH) fue un objetivo de vital importancia. Al entrenar este modelo
con el dataset mencionado y modificado se obtuvieron las m\'etricas expuestas en la tabla 
\ref{tab:comparativa_metricas} y la matriz de confusi\'on de la Figura~\ref{matrix}.

\begin{table}[ht]
    \centering
    \footnotesize
    \caption{M\'etricas obtenidas con el modelo}
    \label{tab:comparativa_metricas}
    \begin{tabularx}{\textwidth}{|>{\raggedright\arraybackslash}X|>{\raggedright\arraybackslash}X|>{\raggedright\arraybackslash}X|>{\raggedright\arraybackslash}X|>{\raggedright\arraybackslash}X|>{\raggedright\arraybackslash}X|>{\raggedright\arraybackslash}X|}
        \hline
        \textbf{Clase} & \textbf{Presici\'on} & \textbf{Sensibilidad} & \textbf{F1} \\
        \hline
        Melanoma & 0.788 & 0.820 & 0.804 \\
        \hline
        Carcinoma Basal & 0.811 & 0.900 & 0.853 \\
        \hline
        Carcinoma Escamoso & 0.833 & 0.781 & 0.806 \\
        \hline
        Otras lesiones & 0.697 & 0.620 & 0.656 \\
        \hline
    \end{tabularx}
\end{table}

\begin{figure}[ht] 
    \centering
    \includegraphics[width=0.6\textwidth]{Graphics/Matrix.png}
    \caption{Matriz de confusi\'on del modelo} 
    \label{matrix}
\end{figure}

Estos resultados se deben a la poca cantidad de im\'agenes del dataset, en la versi\'on original del modelo
se us\'o un dataset de mas de 50000 mil im\'agenes para entrenamiento y pruebas. Las m\'etricas
expresan una sensibilidad no baja al melanoma. Aunque sea menor que la anterior, es un buen n\'umero
cuando se habla en el contexto m\'edico, de cada 100 casos reales de melanoma, el modelo logra 
detectar 82. La categor\'ia de otras lesiones es normal que se mantenga ``baja'', ya que estas
im\'agenes son muy heter\'ogeneas y no siguen siempre un patr\'on visual como puede ser el caso 
de los carcinomas. La matriz de confusi\'on es la explicaci\'on perfecta de lo anterior, ya que 
muchas im\'agenes que deber\'ian ser calificadas como ``otras'' son confundidas con melanomas y 
carcinomas basales.

\subsection{Integraci\'on del Modelo ML en la Aplicaci\'on}
Kaggle, no es solo una plataforma donde extraer y encontrar datasets de inter\'es. En los \'ultimos
a\~nos ha servido de ayuda para desarrolladores e investigadores del campo ``Aprendizaje de M\'aquina'',
permitiendo que estos implementen y entrenen sus modelos utilizando memoria RAM virtual m\'as potente
y navegadores en Internet con velocidades superiores a los 300 MB/s. Para el entrenamiento del modelo
utilizado para DermaUH, se us\'o la platoforma Kaggle como asistencia de GPU e Internet, provocando
que un entrenamiento previsto para d\'ias, se llevara a cabo en cuesti\'on de minutos. 

Una vez finalizado el proceso de entrenamiento, Kaggle genera los archivos correspondientes al 
modelo, los cuales son perfectamente descargables. Una vez descargados, se pueden integrar a 
Django\footnote{Con esta tecnolog\'ia est\'a implementado el backend} utilizando la librer\'ia pytorch.
Por cuestiones de optimizaci\'on, los modelos solamente se cargan una vez iniciada por primera
vez la aplicaci\'on, despu\'es de eso se mantiene activo durante toda la sesi\'on del usuario. 
La Figura~\ref{torchLoadModel} muestra el uso de la librer\'ia torch para consumir los archivos 
del modelo y dejarlos listo para su uso.

\begin{figure}[ht] 
    \centering
    \includegraphics[width=0.6\textwidth]{Graphics/torch.png}
    \caption{Uso de la librer\'ia pytorch para cargar el modelo y dejarlo listo para ejecuci\'on} 
    \label{torchLoadModel}
\end{figure}